% In this file you should put the actual content of the blueprint.
% It will be used both by the web and the print version.
% It should *not* include the \begin{document}
%
% If you want to split the blueprint content into several files then
% the current file can be a simple sequence of \input. Otherwise It
% can start with a \section or \chapter for instance.

\chapter{Compound Protocol Consistency}

We prove that compound protocols (an SWMR global protocol connecting two (SWMR or RCC-Style) cluster protocols
via translation shims) enforce the CMCM (Compound Memory Consistency Model) PPOs:
for any PPO (Preserved Program Order) pair of cache request events at the same
cache with different addresses, if one event is ordered before the other, then
they satisfy a compound linearization order.

% ============================================================
\section{Definitions}
% ============================================================

\begin{definition}[State (Def 5.1)]
  \label{def:State}
  \lean{State}
  \leanok
  (Line 90 in States.lean.)
  A \emph{state} consists of a tuple of a permissions level $p$ (none, read, or write)
  and a coherence flag $c$ (coherent or non-coherent).
  The five states are:
  \begin{itemize}
    \item $\mathrm{SW}$ (Shared-Write): write permissions, coherent,
    \item $\mathrm{MR}$ (Multiple-Read): read permissions, coherent,
    \item $\mathrm{Vd}$ (Valid-dirty): write permissions, non-coherent,
    \item $\mathrm{Vc}$ (Valid-clean): read permissions, non-coherent,
    \item $\mathrm{I}$ (Invalid): no permissions, non-coherent.
  \end{itemize}
\end{definition}

\begin{definition}[Abstract Request (Def 5.2)]
  \label{def:Request}
  \lean{Request}
  \leanok
  (Line 19 in Requests.lean.)
  An \emph{abstract request} consists of:
  \begin{itemize}
    \item A read/write operation,
    \item A coherence flag,
    \item A consistency label (SC, Release, Acquire, or Weak).
  \end{itemize}
  See ValidRequest at line 66, defining a Request and restricting disallowed
  Requests such as a Coherent Acquire or Write Acquire.
\end{definition}

\begin{definition}[Request's Required State (Def 5.3)]
  \label{def:reqHasPerms}
  \lean{Behaviour.reqHasPerms}
  \leanok
  \uses{def:State, def:Request}
  (Line 503 in BehaviourRelationDefs.lean.)
  A request event $e$ \emph{has permissions} if the cache state it is made on
  is at least its Minimum Required State (MRS).
  Let's start by defining a simple Required State (RS) of a request `r'.
  Let the read/write operation `rw' and coherence flag `c' of `r' be the State tuple (`rw', `c')
  to define its RS.

  For coherent requests (ex. a SC Write from a SWMR protocol, or Coherent Release Writes from RCC-O)
  we use the simple RS definition to define the MRS.
  This is reflected in the first case in the inductive `Behaviour.reqHasPerms'.

  For non-coherent acquire/release/weak writes from RCC protocols, the request's MRS
  additionally requires the state is coherent.
  This is reflected in the second case in the inductive `Behaviour.reqHasPerms'.

  For non-coherent weak read/write requests uses the simple RS definition and specifies
  the cache state is Read-Non-Coherent. The state is not Vd.
  This is reflected in the third case in the inductive `Behaviour.reqHasPerms'.
\end{definition}

\begin{definition}[Local Linearization Point (Def 5.4)]
  \label{def:linearizationEventOfRequest}
  \lean{Behaviour.linearizationEventOfRequest}
  \leanok
  \uses{def:reqHasPerms}
  (Line 1083 in BehaviourRelationProofs.lean.)
  The \emph{linearization event} of a request event $e_{\mathrm{req}}$.
  Either the cluster cache request event linearizes in the cache (on a state with coherent permissions),
  or at the directory (without coherent permissions). Linearizing at the directory means
  either linearizing before, during, or after the request (as per inductive `Behaviour.dirAccessOfRequest'
  on line 569 of BehaviourRelationDefs.lean). More succintly in `Behaviour.linearizationEventOfRequest':
  \begin{itemize}
    \item If $e_{\mathrm{req}}$ is made on a coherent state with sufficient permissions,
      the linearization event is $e_{\mathrm{req}}$ itself (\texttt{requestLin}).
    \item Otherwise, the linearization event is the directory event $e_{\mathrm{dir}}$
      corresponding to $e_{\mathrm{req}}$ (\texttt{dirLin}).
  \end{itemize}
\end{definition}

\begin{definition}[Global Linearization Point (Def 5.5)]
  \label{def:ClusterRequestLinearizationEvent}
  \lean{ClusterRequestLinearizationEvent}
  \leanok
  \uses{def:linearizationEventOfRequest, def:ShimClusterToGlobal, def:CompoundProtocol}
  (Line 139 in CompoundLinearization.lean.)
  A cluster request's cluster linearization event classifies where a cluster cache request
  globally linearizes:
  \begin{itemize}
    \item \textbf{At the Cluster Cache} (\texttt{clusterCacheLin}): the request has permissions and
      linearizes locally; the linearization event is the cache event itself.
    \item \textbf{At the Cluster Directory or Beyond} (\texttt{clusterDirLin}): the request does not
      have permissions locally and must go through the directory (and possibly
      the global protocol). The request then either globally linearizes at the directory
      (with global permissions already) or linearizes at the global directory (needs
      to access the global directory for permissions).
  \end{itemize}
  In the case where the cluster request linearizes at the cluster directory or global protocol
  is defined by the inductive `CompoundProtocol.clusterDirectoryLinearizationEvent' at line 98
  in the same file.
\end{definition}

% ============================================================
\section{Well-Formedness Axioms}
% ============================================================

\begin{definition}[Ordered Directory Events (Well-Formedness Axiom 1)]
  \label{def:DirectoryEventAreOrdered}
  \lean{DirectoryEvent.AreOrdered}
  \leanok
  (Line 305 in EventRelations.lean.)
  Events at the same directory address are totally ordered:
  for any two directory events $de_1, de_2$ at the same address,
  either $de_1$ is ordered before $de_2$ or vice versa.
  This is stated by the structure fields in `DirectoryEvent.AreOrdered',
  `sameDirectoryEntry' and `ordered'.
\end{definition}

\begin{definition}[Ordered/Encapsulated Cache Events (Well-Formedness Axiom 2)]
  \label{def:CacheEventAreOrdered}
  \lean{CacheEvent.AreOrdered}
  \leanok
  (Line 341 in EventRelations.lean.)
  For any two cache events $e_1, e_2$ at the same cache entry,
  either $e_1$ is encapsulated by $e_2$, $e_2$ is encapsulated by
  $e_1$, $e_1$ is ordered before $e_2$, or $e_2$ is ordered before $e_1$.
  This is stated by the structure fields in `CacheEvent.AreOrdered',
  `sameCacheEntry' and `ordered'.
\end{definition}

\begin{definition}[Cache Events Encapsulate Corresponding Directory Event (Well-Formedness Axiom 3)]
  \label{def:requestDirectoryEvent}
  \lean{Behaviour.requestDirectoryEvent}
  \leanok
  \uses{def:DirectoryEventAreOrdered, def:CacheEventAreOrdered}
  (Line 55 in BehaviourRelationDefs.lean.)
  A cache operation event encapsulates a corresponding directory event,
  whose request field, address, downgrade flag, and protocol match
  the cache event.
  This is reflected in the fields of structure `Behaviour.requestDirectoryEvent'.
\end{definition}

\begin{definition}[Cache Events Access Dir If Insufficient Permissions (Well-Formedness Axiom 4)]
  \label{def:axRequestAccessesDirectory}
  \lean{Behaviour.axRequestAccessesDirectory}
  \leanok
  \uses{def:reqHasPerms, def:requestDirectoryEvent}
  (Line 225 in BehaviourRelationDefs.lean.)
  If a cache event does not have sufficient permissions (its state is below
  the Minimum Required State), it accesses the directory.
\end{definition}

\begin{definition}[Cache Events of Directory Responses Arrive In-Order (Well-Formedness Axiom 5)]
  \label{def:axiom5_dir_responses_ordered}
  This axiom states that cache events of directory responses arrive in-order.
  It does not have a separate Lean declaration; it is coincidentally
  implemented by Axiom 1 (\ref{def:DirectoryEventAreOrdered}) combined with
  the Downgrade Abstract State Invariants --- downgrades are encapsulated by
  a directory event, so they are in order because directory events are ordered.
\end{definition}

\begin{definition}[Directory Responses to Caches Are Ordered (Well-Formedness Axiom 6)]
  \label{def:encapGrantAfterDirEvent}
  \lean{Event.encapGrantAfterDirEvent}
  \leanok
  \uses{def:DirectoryEventAreOrdered, def:axiom5_dir_responses_ordered}
  (Line 532 in EventRelations.lean.)
  A request event encapsulates a grant event that is ordered after
  the directory event, and the grant ends the request.
\end{definition}

\begin{definition}[Cache with Perms Has Predecessor That Obtained Perms (Well-Formedness Axiom 7)]
  \label{def:exists_predecessor_setting_state}
  \lean{Behaviour.exists_predecessor_setting_state}
  \leanok
  \uses{def:reqHasPerms, def:axRequestAccessesDirectory}
  (Line 494 in BehaviourRelationProofs.lean.)
  If a non-downgrade request event has permissions, there exists an
  immediate predecessor that obtained those permissions by accessing
  the directory, and all intermediate events maintain sufficient state.
\end{definition}

\begin{definition}[Weak Request on Vd Has Successor WB or Get SW (Well-Formedness Axiom 8)]
  \label{def:exists_vd_successor_wb_or_get_sw}
  \lean{Behaviour.exists_vd_successor_wb_or_get_sw}
  \leanok
  (Line 504 in BehaviourRelationProofs.lean.)
  A weak request on Write-Non-Coherent state ordered before a write-back evict
  linearizes at the write-back evict.
  There exists an immediate bottom successor that encapsulates a corresponding
  directory event.
\end{definition}

\begin{definition}[Acquire Invalidation OB Weak Read Predecessor (Well-Formedness Axiom 9)]
  \label{def:acquire_invalidation_ob_weak_read}
  \lean{CompoundProtocol.acquire_invalidation_OrderedBefore_weak_read_predecessor_that_gets_perms}
  \leanok
  (Line 8 in Protocol.lean.)
  A non-coherent weak read ordered after a cache entry invalidation (from an
  acquire) will linearize after the invalidation: the invalidation event is
  ordered before the predecessor that gets permissions.
\end{definition}

% ============================================================
\section{Abstract State Invariants}
% ============================================================

\begin{definition}[NC Weak Write Linearizes After Cache Event (Abstract State Invariant 1)]
  \label{def:vdCacheEntryWBOrGetSWLaterWrapper}
  \lean{Behaviour.vdCacheEntryWBOrGetSWLaterWrapper}
  \leanok
  (Line 238 in BehaviourRelationDefs.lean.)
  Non-coherent weak write request linearizes after its cache event.
  A cache entry in Vd state has an immediate bottom successor that is
  either a VdWriteBack or a CoherentWrite (which would transition the state).
\end{definition}

\begin{definition}[Coherent Write Dir Events Result in Downgrades (Abstract State Invariant 2)]
  \label{def:coherentWriteAtDirectoryEncapDowngrades}
  \lean{Behaviour.coherentWriteAtDirectoryEncapDowngrades}
  \leanok
  \uses{def:requestDirectoryEvent}
  (Line 294 in BehaviourRelationDefs.lean.)
  Coherent write directory events result in downgrades.
  A coherent write request to the directory results in downgrades at other
  caches. On SW state, it forwards to the owner; on MR state, it forwards
  to all sharers.
\end{definition}

\begin{definition}[Coherent Read Dir on SW Downgrades Owner (Abstract State Invariant 3)]
  \label{def:coherentReadDirDowngradeOthers}
  \lean{Behaviour.coherentReadDirDowngradeOthers}
  \leanok
  \uses{def:requestDirectoryEvent}
  (Line 324 in BehaviourRelationDefs.lean.)
  Coherent read directory event on Write-Coherent state downgrades the owner.
  A coherent read request to the directory on SW state results in
  a downgrade forwarded to the owner.
\end{definition}

\begin{definition}[NC Requests on SW Downgrade Owner (Abstract State Invariant 4)]
  \label{def:nonCoherentRequestDowngradeOthers}
  \lean{Behaviour.nonCoherentRequestDowngradeOthers}
  \leanok
  \uses{def:requestDirectoryEvent}
  (Line 345 in BehaviourRelationDefs.lean.)
  Non-coherent requests on Write-Coherent state will downgrade the owner.
  Non-coherent write/read requests on SW directory state result in
  downgrades forwarded to the previous owner.
\end{definition}

\begin{definition}[Protocol Enforces SWMR (Abstract State Invariant 5)]
  \label{def:SWMR_wrapper}
  \lean{SWMR.wrapper}
  \leanok
  \uses{def:State}
  (Line 96 in SWMR.lean.)
  The protocol enforces Single-Writer Multiple-Reader (SWMR):
  at any point in time, cache states across the protocol instance
  cannot simultaneously have both SW and MR, and at most one cache
  can be in SW state.
\end{definition}

% ============================================================
\section{Consistency Label Axioms}
% ============================================================

\begin{definition}[Acquire Invalidates Other Entries (Consistency Label Axiom 1)]
  \label{def:acqInvalOtherEntries}
  \lean{Behaviour.acqInvalOtherEntries}
  \leanok
  \uses{def:requestDirectoryEvent, def:CacheEventAreOrdered}
  (Line 374 in BehaviourRelationDefs.lean.)
  Acquire cache operations access the directory, then invalidate
  Read-Non-Coherent cache entries.
  An acquire cache operation invalidates Vc cache entries
  at other addresses after accessing the directory: for every other address,
  there is a Vc invalidation event encapsulated by the acquire, ordered
  after the directory event.
\end{definition}

\begin{definition}[NC Release Writes Back Other Entries (Consistency Label Axiom 2)]
  \label{def:relWriteBackOtherEntries}
  \lean{Behaviour.relWriteBackOtherEntries}
  \leanok
  \uses{def:requestDirectoryEvent}
  (Line 385 in BehaviourRelationDefs.lean.)
  Non-coherent release cache operations write back Write-Non-Coherent entries,
  then access the directory.
  A non-coherent release cache operation writes back Vd
  entries at other addresses before accessing the directory: for every other
  address, there is a Vd write-back event encapsulated by the release, ordered
  before the directory event.
\end{definition}

\begin{definition}[L-RCC Coherent Release Downgrade Write-Back (Consistency Label Axiom 3)]
  \label{def:coherentRelDowngradeWriteBackOthers}
  \lean{Behaviour.coherentRelDowngradeWriteBackOthers}
  \leanok
  (Line 407 in BehaviourRelationDefs.lean.)
  L-RCC variant with Coherent Release Write and Non-Coherent Weak Write.
  When receiving a downgrade, a coherent release writes back Vd entries
  at other addresses. This is the L-RCC variant where the protocol interface
  contains both a coherent release and a non-coherent weak write.
\end{definition}

% ============================================================
\section{Shim Axioms}
% ============================================================

\begin{definition}[Cluster to Global Propagation Shim (Shim Axiom 1)]
  \label{def:ShimClusterToGlobal}
  \lean{Behaviour.Shim.ClusterToGlobal}
  \leanok
  \uses{def:CompoundProtocol}
  (Line 128 in BehaviourShim.lean.)
  Cluster directory events are translated to request events at the
  corresponding cache in the global protocol. Two cases:
  \begin{itemize}
    \item \textbf{encapGlobalCache}: the global cache lacks permissions,
      so the directory event encapsulates a global cache access.
    \item \textbf{noGlobalCache}: the global cache already has permissions,
      and a prior global cache event with permissions exists.
  \end{itemize}
\end{definition}

\begin{definition}[Global to Cluster Propagation Shim (Shim Axiom 2)]
  \label{def:ShimGlobalToCluster}
  \lean{Behaviour.Shim.GlobalToCluster}
  \leanok
  \uses{def:CompoundProtocol}
  (Line 903 in BehaviourShim.lean.)
  Downgrades at a global cache are translated to cluster directory accesses.
  The translation depends on whether the cluster protocol has both coherent
  write and read, or only non-coherent reads.
\end{definition}

\begin{definition}[Ordered Before]
  \label{def:OrderedBefore}
  \lean{Event.OrderedBefore}
  \leanok
  An event $e_1$ is \emph{ordered before} $e_2$ if $e_1$ completes before $e_2$
  starts: $e_1.\mathrm{oEnd} < e_2.\mathrm{oStart}$.
\end{definition}

\begin{definition}[Encapsulates]
  \label{def:Encapsulates}
  \lean{Event.Encapsulates}
  \leanok
  An event $e_1$ \emph{encapsulates} $e_2$ if $e_1$ starts before $e_2$ starts and
  $e_2$ ends before $e_1$ ends:
  $e_1.\mathrm{oStart} < e_2.\mathrm{oStart} \;\wedge\; e_2.\mathrm{oEnd} < e_1.\mathrm{oEnd}$.
\end{definition}

\begin{definition}[Compound Protocol]
  \label{def:CompoundProtocol}
  \lean{CompoundProtocol}
  \leanok
  A \emph{compound protocol} bundles:
  \begin{itemize}
    \item A global protocol and two cluster protocols,
    \item Shim axioms for translating between cluster and global protocols,
    \item A linearization function assigning each request event a linearization event,
    \item A compound linearization event function,
    \item Axioms: the global protocol is SWMR,
      acquire events on a shared-write cannot linearize at cache,
      and weak write / non-coherent release pairs cannot both linearize at cache.
  \end{itemize}
\end{definition}

\begin{definition}[Lazy Compound Linearization Order]
  \label{def:lazyCompoundLinearizationOrder}
  \lean{CompoundProtocol.lazyCompoundLinearizationOrder}
  \leanok
  \uses{def:CompoundProtocol}
  A \emph{lazy compound linearization order} between a linearization event $e_{\mathrm{lin}_1}$
  and a second event $e_2$ holds if: for every immediate bottom predecessor $e_3$ of
  $e_2$'s directory event at the same address, $e_{\mathrm{lin}_1}$ finishes before $e_3$'s
  directory event.
\end{definition}

\begin{definition}[Compound Linearization Order]
  \label{def:CompoundLinearizationOrder}
  \lean{CompoundProtocol.CompoundLinearizationOrder}
  \leanok
  \uses{def:lazyCompoundLinearizationOrder, def:ClusterRequestLinearizationEvent}
  For a PPO pair $(e_1, e_2)$, \emph{compound linearization order} holds if
  either their compound linearization events are ordered
  ($e_{\mathrm{lin}_1}$ is ordered before $e_{\mathrm{lin}_2}$),
  or the lazy compound linearization order holds between $e_{\mathrm{lin}_1}$ and $e_2$.
\end{definition}

% ============================================================
\section{Auxiliary Structures}
% ============================================================

\begin{definition}[Coherent Write Succeeded by Directory Request]
  \label{def:coherentWriteSucceededByDirectoryRequest}
  \lean{Behaviour.coherentWriteSucceededByDirectoryRequest}
  \leanok
  Structure capturing the auxiliary data for a coherent write event
  that is succeeded by a directory request.
\end{definition}

\begin{definition}[RCCO-Style Downgrade]
  \label{def:rccOStyleDowngrade}
  \lean{Event.rccOStyleDowngrade}
  \leanok
  Structure capturing the chain of events in an RCCO-style downgrade:
  a directory event, a downgrade event, and a write-back event.
\end{definition}

\begin{definition}[Write-Back at Event]
  \label{def:writeBackAtEvent}
  \lean{Event.writeBackAtEvent}
  \leanok
  Structure recording a VdWriteBack at a given event's entry.
\end{definition}

% ============================================================
\section{Contradiction Lemmas}
% ============================================================

\begin{lemma}[Contradiction: Has Perms and No Perms]
  \label{lem:contradiction_has_perms_and_no_perms}
  \lean{Event.contradiction_of_has_perms_and_no_perms}
  \leanok
  A coherent event cannot simultaneously have and lack permissions.
\end{lemma}

\begin{lemma}[Contradiction: NC Release Has Perms and No Perms]
  \label{lem:contradiction_nc_release_perms}
  \lean{Event.contradiction_of_nc_release_request_has_perms_and_no_perms}
  \leanok
  \uses{lem:contradiction_has_perms_and_no_perms}
  A non-coherent release request cannot simultaneously
  have and lack permissions.
\end{lemma}

\begin{lemma}[Contradiction: Weak Write Has Perms and No Perms]
  \label{lem:contradiction_weak_write_perms}
  \lean{Event.contradiction_of_weak_write_request_has_perms_and_no_perms}
  \leanok
  \uses{lem:contradiction_has_perms_and_no_perms}
  A weak write request cannot simultaneously have and lack permissions.
\end{lemma}

\begin{lemma}[Contradiction: Acquire Has Perms and No Perms]
  \label{lem:contradiction_acquire_perms}
  \lean{Event.contradiction_of_acquire_request_has_perms_and_no_perms}
  \leanok
  \uses{lem:contradiction_has_perms_and_no_perms}
  An acquire request cannot simultaneously have and lack permissions.
\end{lemma}

\begin{lemma}[Contradiction: Coherent Has Perms and No Perms]
  \label{lem:contradiction_coherent_perms}
  \lean{Event.contradiction_of_coherent_request_has_perms_and_no_perms}
  \leanok
  \uses{lem:contradiction_has_perms_and_no_perms}
  A coherent request cannot simultaneously have and lack permissions.
\end{lemma}

% ============================================================
\section{Encapsulation Lemmas}
% ============================================================

\begin{lemma}[Directory Lin Encapsulates Global Directory Lin]
  \label{lem:dir_lin_encap_global_dir_lin}
  \lean{CompoundProtocol.dir_lin_encap_global_dir_lin}
  \leanok
  \uses{def:Encapsulates, def:ClusterRequestLinearizationEvent}
  If a cluster directory event encapsulates the global directory
  linearization event, then the encapsulation relation is preserved
  through the linearization chain.
\end{lemma}

\begin{lemma}[Event Encaps Dir Lin Encaps Global Dir Lin]
  \label{lem:e_encap_dir_lin_encap_global_dir_lin}
  \lean{CompoundProtocol.e_encap_dir_lin_encap_global_dir_lin}
  \leanok
  \uses{lem:dir_lin_encap_global_dir_lin}
  If a cache event encapsulates a cluster directory event, and the cluster
  directory event encapsulates the global directory linearization event,
  then the cache event encapsulates the global directory linearization event.
\end{lemma}

\begin{lemma}[Request Encapsulates Compound Linearization Event]
  \label{lem:request_encapsulates_cmp_lin}
  \lean{CompoundProtocol.request_encapsulates_compound_linearization_event}
  \leanok
  \uses{lem:e_encap_dir_lin_encap_global_dir_lin}
  A request event encapsulates its compound linearization event
  (generic version for all request types).
\end{lemma}

\begin{lemma}[NC Release Encapsulates Compound Lin Event]
  \label{lem:nc_release_encapsulates_cmp_lin}
  \lean{CompoundProtocol.nc_release_request_encapsulates_compound_linearization_event}
  \leanok
  \uses{lem:e_encap_dir_lin_encap_global_dir_lin, lem:contradiction_nc_release_perms}
  A non-coherent release request encapsulates its compound linearization event.
\end{lemma}

\begin{lemma}[Acquire Encapsulates Compound Lin Event]
  \label{lem:acquire_encapsulates_cmp_lin}
  \lean{CompoundProtocol.acquire_request_encapsulates_compound_linearization_event}
  \leanok
  \uses{lem:e_encap_dir_lin_encap_global_dir_lin, lem:contradiction_acquire_perms}
  An acquire request encapsulates its compound linearization event.
\end{lemma}

\begin{lemma}[Coherent Encapsulates Compound Lin Event]
  \label{lem:coherent_encapsulates_cmp_lin}
  \lean{CompoundProtocol.coherent_request_encapsulates_compound_linearization_event}
  \leanok
  \uses{lem:e_encap_dir_lin_encap_global_dir_lin, lem:contradiction_coherent_perms}
  A coherent request encapsulates its compound linearization event.
\end{lemma}

% ============================================================
\section{Linearization Ordering Lemmas}
% ============================================================

\begin{lemma}[Ordered Before and Linearizes at Cache]
  \label{lem:ob_linearizes_at_cache}
  \lean{CompoundProtocol.compound_linearization_order_of_events_ordered_before_and_linearizes_at_cache}
  \leanok
  \uses{def:CompoundLinearizationOrder, def:OrderedBefore}
  If both events linearize at cache and $e_1$ is ordered before $e_2$,
  then compound linearization order holds.
  The linearization events are the events themselves.
\end{lemma}

\begin{lemma}[Cluster Dir Encap Global Lin at Dir]
  \label{lem:cdir_encap_glin_at_dir}
  \lean{CompoundProtocol.cdir_encap_glin_of_cdir_linearize_at_dir}
  \leanok
  \uses{lem:dir_lin_encap_global_dir_lin}
  A cluster directory event encapsulates its compound linearization event
  when linearizing at directory.
\end{lemma}

\begin{lemma}[Weak Write/Read and NC Release Linearize at Directory]
  \label{lem:weak_and_nc_release_at_dir}
  \lean{CompoundProtocol.weak_write_or_read_and_nc_release_linearize_at_directory}
  \leanok
  \uses{lem:cdir_encap_glin_at_dir}
  If a weak write or read and a non-coherent release both linearize at
  directory, their linearization events are ordered.
\end{lemma}

\begin{lemma}[Acquire and Weak Request Linearize at Directory]
  \label{lem:acquire_and_weak_at_dir}
  \lean{CompoundProtocol.acquire_and_weak_request_linearize_at_directory}
  \leanok
  \uses{lem:cdir_encap_glin_at_dir}
  If an acquire and a weak request both linearize at directory,
  their linearization events are ordered.
\end{lemma}

\begin{lemma}[Weak Write Ordered Before VdWriteBack]
  \label{lem:weak_write_ob_vd_wb}
  \lean{CompoundProtocol.weak_write_OrderedBefore_vd_write_back'}
  \leanok
  A weak write that is ordered before a VdWriteBack linearizes before the
  write-back.
\end{lemma}

\begin{lemma}[State Before Vd of No Between Events]
  \label{lem:state_before_vd_no_btn}
  \lean{CompoundProtocol.state_before_Vd_of_no_btn_events_on_vd_produce_vd}
  \leanok
  If there are no between-events on a Vd entry, the state before the Vd
  event produces a Vd state.
\end{lemma}

\begin{lemma}[Lazy Coherent Release Ordering (Helper)]
  \label{lem:lazy_coherent_release_ordering_helper}
  \lean{CompoundProtocol.lazy_coherent_release_ordering'}
  \leanok
  Helper lemma for lazy coherent release ordering.
\end{lemma}

\begin{lemma}[Lazy Coherent Release Ordering]
  \label{lem:lazy_coherent_release_ordering}
  \lean{CompoundProtocol.lazy_coherent_release_ordering}
  \leanok
  \uses{lem:lazy_coherent_release_ordering_helper}
  Establishes the lazy linearization ordering for coherent releases.
\end{lemma}

\begin{lemma}[Weak Request to Dir and Coherent Release at Cache]
  \label{lem:weak_req_dir_crel_cache}
  \lean{CompoundProtocol.weak_request_to_directory_and_coherent_release_at_cache}
  \leanok
  When a weak request linearizes at directory and a coherent release
  linearizes at cache, compound linearization order holds.
\end{lemma}

\begin{lemma}[Contradiction: Weak Write Encapsulating Directory Event]
  \label{lem:contradiction_weak_write_encap_dir}
  \lean{CompoundProtocol.contradiction_of_weak_write_encapsulating_corresponding_directory_event}
  \leanok
  A contradiction arises if a weak write encapsulates its corresponding
  directory event.
\end{lemma}

% ============================================================
\section{Per-Case Compound Linearization Lemmas}
% ============================================================

\begin{lemma}[CLO: Two Events Encapsulating Their Lin Events]
  \label{lem:clo_two_events_encap_lin}
  \lean{CompoundProtocol.CompoundLinearizationOrder_of_two_events_that_encapsulate_their_cmp_linearization_event}
  \leanok
  \uses{lem:ob_linearizes_at_cache, lem:nc_release_encapsulates_cmp_lin,
        lem:acquire_encapsulates_cmp_lin, lem:coherent_encapsulates_cmp_lin}
  Handles PPO cases where both events are non-coherent release, acquire,
  or coherent (e.g.\ SC$\times$SC, Rel$\times$Acq, Acq$\times$Rel).
  Shows that each event encapsulates its compound linearization event,
  so ordered-before is preserved to the linearization events.
\end{lemma}

\begin{lemma}[CLO: Weak Read and NC Release]
  \label{lem:clo_weak_read_nc_release}
  \lean{CompoundProtocol.CompoundLinearizationOrder_of_weak_read_and_non_coherent_release}
  \leanok
  \uses{lem:ob_linearizes_at_cache, lem:weak_and_nc_release_at_dir}
  Weak read ordered before a non-coherent release satisfies compound
  linearization order.
\end{lemma}

\begin{lemma}[CLO: Weak Write or Read and NC Release]
  \label{lem:clo_weak_wr_nc_release}
  \lean{CompoundProtocol.CompoundLinearizationOrder_of_weak_write_or_read_and_non_coherent_release}
  \leanok
  \uses{lem:ob_linearizes_at_cache, lem:weak_and_nc_release_at_dir, lem:clo_weak_read_nc_release}
  Generic lemma: weak write or read before a non-coherent release.
  Handles all sub-cases (cache$\times$cache, dir$\times$cache,
  cache$\times$dir, dir$\times$dir linearization).
\end{lemma}

\begin{lemma}[CLO: Weak Request and NC Release]
  \label{lem:clo_weak_req_nc_release}
  \lean{CompoundProtocol.CompoundLinearizationOrder_of_weak_request_and_non_coherent_release}
  \leanok
  \uses{lem:clo_weak_wr_nc_release}
  Weak non-coherent request ordered before a non-coherent release satisfies
  compound linearization order.
\end{lemma}

\begin{lemma}[CLO: Acquire and Weak Request]
  \label{lem:clo_acquire_weak}
  \lean{CompoundProtocol.CompoundLinearizationOrder_of_acquire_and_weak_request}
  \leanok
  \uses{lem:ob_linearizes_at_cache, lem:acquire_and_weak_at_dir,
        lem:contradiction_weak_write_encap_dir}
  Acquire ordered before a weak non-coherent request satisfies compound
  linearization order.
\end{lemma}

\begin{lemma}[CLO: Weak Request and Coherent Release]
  \label{lem:clo_weak_req_crel}
  \lean{CompoundProtocol.CompoundLinearizationOrder_of_weak_request_and_coherent_release}
  \leanok
  \uses{lem:ob_linearizes_at_cache, lem:weak_req_dir_crel_cache,
        lem:lazy_coherent_release_ordering}
  Weak non-coherent request ordered before a coherent release satisfies
  compound linearization order.
\end{lemma}

% ============================================================
\section{Main Results}
% ============================================================

\begin{lemma}[PPO Cluster Events Satisfy Compound Linearization]
  \label{lem:ppo_cluster_events}
  \lean{CompoundProtocol.ppo_cluster_events_satisfy_CompoundLinearizationOrder}
  \leanok
  \uses{def:CompoundLinearizationOrder,
        lem:clo_two_events_encap_lin,
        lem:clo_weak_req_nc_release,
        lem:clo_weak_req_crel,
        lem:clo_acquire_weak}
  For same-cluster PPO cache events satisfying ordered-before,
  the compound linearization order holds.
  The proof is by exhaustive case match on the request types of $e_1$ and $e_2$,
  dispatching to the appropriate per-case lemma.
\end{lemma}

\begin{theorem}[Compound Protocol Enforces Consistency (Theorem 6.1)]
  \label{thm:enforce_compound_consistency}
  \lean{CompoundProtocol.enforce_compound_consistency}
  \leanok
  \uses{lem:ppo_cluster_events}
  (Line 8 in CompositionalMCM.lean.)
  For any PPO pair of non-downgrade cache events $e_1, e_2$ at the same cache
  with different addresses in a behaviour $b$ of a compound protocol:
  if $e_1$ is ordered before $e_2$, then compound linearization order holds
  between $e_1$ and $e_2$.
\end{theorem}
